\subsection{Домашние задачи на девятую неделю.}
\begin{problem}
    Используя неравенство Маркова, вывести показательную форму неравенства Чебышёва:
    \[
        \P(X \geq a) \leq \dfrac{\Ex e^{tX}}{e^{ta}}, t > 0, a \in \R.
    \]
\end{problem}
\begin{solution}
    Заметим, что в силу строгой монотонности экспоненты следует $X \geq a \Longleftrightarrow e^{tX} \geq e^{ta}$ при положительных $t$.
    Из этого следует
    \[
        \P(X \geq a) = \P(e^{tX} \geq e^{ta}).
    \]
    Применим неравенство Маркова и получим утверждение задачи:
    \[
        \P(X \geq a) = \P(e^{tX} \geq e^{ta}) \leq \dfrac{\Ex e^{tX}}{e^{ta}}.
    \]
\end{solution}
\begin{problem}
    Оказывается, что ЗБЧ может быть выполнен и в случае, когда $X_n$ не обязательно независимы и одинаково распределены.
    Доказать, что если $\cov(X_i, X_j) \leq 0$ при $i \neq j$ и $\Ex X_n \ra m$, $\sigma^2_{n} = \D X_n = o(n)$, то $S_n/n \ra m$ по вероятности.
\end{problem}
\begin{solution}
    Для начала, заметим, что так как $\Ex(X_n) \ra m, n \ra +\infty$, то и $\Ex(S_n)/n \ra m, n \ra +\infty$ -- поскольку если у последовательности существует предел в обычном смысле, то он существует и в смысле Чезаро и они равны. \\
    Пусть $Y_n = S_n/n - \Ex(S_n)/n$.
    Зафиксируем произвольное $\epsilon > 0$.
    Согласно неравенству Чебышёва:
    \[
        \P(|Y_n| \geq \epsilon) \leq \dfrac{\D Y_n}{\epsilon^2}.
    \]
    Заметим, что $\Ex Y_n = 0$.
    Распишем дисперсию $Y_n$:
    \[
        DY_n = \Ex Y_n^2 = \Ex \biggr( \sum\limits_{k = 1}^n \dfrac{\overset{\circ}{X_n}}{n} \biggr)^2 = \dfrac{1}{n^2}\biggr(\sum\limits_{k = 1}^n \sigma_k^2 + \sum\limits_{i, j = 1, 1; i \neq j}^{n, n} \cov (X_i, X_j) \biggr) \leq \dfrac{1}{n^2}\sum\limits_{k = 1}^n \sigma^2_k \ra 0, n \ra +\infty.
    \]
    Переход в неравенстве сделан в силу того, что $X_i$ и $X_j$ имеют отрицательную корреляцию при $i \neq j$; предельный переход корректен, поскольку $\sum\limits_{k = 1}^n \sigma^2_k = o(n^2)$ -- а снизу у нас как раз $n^2$.
    Таким образом $Y_n \xrightarrow{\P} 0$, а следовательно, с учётом того, что $\Ex S_n/n \ra m$, то и $S_n/n \xrightarrow{\P} m$ -- что и требовалось доказать.
\end{solution}
\begin{problem}
    Пусть $X_1, X_2, \ldots$~---~независимые случайные величины и $X_n \ra X$ по вероятности.
    Доказать, что $X = const$ п.н.
\end{problem}
\begin{solution}
    Является следствием закона <<0-1>>.
\end{solution}
\begin{problem}
    Рассмотрим случайные величины $X_n \sim U(-n, n)$.
    Найти их характеристические функции $\phi_{X_n}$.
    Сходятся ли эти случайные величины по распределению?
\end{problem}
\begin{solution}
    Вычислим $\phi_n(t) = \Ex e^{it X_n}$:
    \[
        \phi_n(t) = \int_{-n}^{n} \dfrac{e^{itx}}{2n}dx = \dfrac{1}{2int}\big(e^{int} - e^{-int}\big) = \dfrac{\sin nt}{nt}.
    \]
    Если $t \neq 0$, то $\phi_n(t) \ra 0, n \ra  +\infty$; если же $t = 0$, то $\phi_n(0) = 1$ и $\phi_n(0) \ra 1, n \ra +\infty$.
    Таким образом $\phi_n \ra I_{\{0\}}, n \ra 0$.
    По теореме Бохнера-Хинчина $I_{\{0\}}$ не является характеристической функцией распределения (поскольку нет непрерывности в окрестности нуля), а следовательно $\phi_n$ не сходятся по распределению.
\end{solution}
\begin{problem}
    Доказать, что если $X_n$ сходится по распределению к константе, то $X_n$ сходится к ней и по вероятности.
\end{problem}
\begin{solution}
    Пусть $X_n \xrightarrow{d} c$.
    По определению $F_{X_n} \ra F_c, n \ra +\infty$ во всех точках непрерывности.
    Поскольку $F_c = I_{[c, +\infty)}$, то $F_c$ непрерывна всюду, кроме $c$.
    Таким образом:
    \[
        F_{X_n}(t) \ra \begin{cases} 0, & t < c \\
                           1, & t > c.
        \end{cases}, n \ra +\infty.
    \]
    Поскольку
    \[
        \P(|X_n - c| \geq \epsilon) = \P(X_n \leq c - \epsilon) + \P(X_n \geq c + \epsilon).
    \]
    И
    \begin{multline*}\P(X_n \leq c - \epsilon) = \int_{-\infty}^{c - \epsilon}dF_{X_n} = F_{X_n}(c - \epsilon) - F_{X_n}(-\infty) = F_{X_n}(c - \epsilon), \\ \P(X_n \geq c + \epsilon) = \int_{c + \epsilon}^{+\infty} dF_{X_n} = F_{X_n}(+\infty) - F_{X_n}(c + \epsilon) = 1 - F_{X_n}(c + \epsilon).
    \end{multline*}
    То в силу уже отмеченных поточечных сходимостей мы получаем сходимость по вероятности для $X_n$ к $c$.
\end{solution}
\begin{problem}
    Исследуйте на сходимость в среднеквадратичном среднего геометрического независимых случайных величин $\{X_n\}_{n \in \N}$, каждая из которых $\sim U(0, 1)$.
\end{problem}
\begin{solution}
    Средним геометрическим $X_1, \ldots, X_n$ является:
    \[
        \sqrt[n]{\prod_{k = 1}^n X_k} = \exp\biggr(-{\dfrac{1}{n}\sum\limits_{k = 1}^n -\ln X_k}\biggr).
    \]
    Нетрудно заметить, что $Y_k = -\ln X_k \sim Exp(1)$.
    Поскольку $Y_k$ все -- независимо одинаково распределённые случайные величины с $\Ex Y_k = 1$, то по УЗБЧ в форме Колмогорова
    \[
        \dfrac{1}{n}\sum\limits_{k = 1}^n Y_k \xrightarrow{a.s.} 1, n \ra +\infty.
    \]
    Из непрерывности экспоненты следует
    \[
        \exp\biggr(-\dfrac{1}{n}\sum\limits_{k = 1}^n Y_k\biggr) \xrightarrow{a.s.} e^{-1}, n \ra +\infty.
    \]
    Пусть $Z_n = \biggr(\exp\biggr(-\dfrac{1}{n}\sum\limits_{k = 1}^n Y_k\biggr) - e^{-1}\biggr)^2$; $Z_n \xrightarrow{a.s.} 0, n \ra +\infty$. \\
    Поскольку
    \[
        \int Z_{n}d\P \leq \int \exp \biggr(-\dfrac{2}{n}\sum\limits_{k = 1}^n Y_k\biggr)d\P + 2e^{-1}\int \exp \biggr(-\dfrac{1}{n}\sum\limits_{k = 1}^n Y_k\biggr)d\P + e^{-2} < +\infty.
    \]
    То у нас есть интегрируемая мажоранта и по теореме Лебега $\int Z_n \ra 0, n \ra +\infty$ -- что и показывает сходимость по $L_2$-норме для среднего геометрического равномерно распределённых случайных величин.
\end{solution}
\begin{note}
    Мне не нравится доказательство того, что здесь из a.s. следует $L_2$; кажется, это не совсем верно. Надо переделать.
\end{note}
\begin{problem}
    Следует ли из сходимости случайных величин п.н. сходимость их математических ожиданий?
    А из сходимости случайных величин по вероятности?
    Считайте, что все рассматриваемые случайные величины последовательности имеют конечное математическое ожидание.
\end{problem}
\begin{solution} \leavevmode
    \begin{itemize}
        \item Покажем, что $X_n \xrightarrow{a.s.} X, n \ra +\infty \not\Ra X_n \xrightarrow{L_1} X, n \ra +\infty$.
        Пусть $X_n = n I_{[0, 1/n)}$.
        Тогда $X_n \xrightarrow{a.s.} 0, n \ra +\infty$, а значит если $X_n \xrightarrow{L_1} X$, то $X = 0$ (п.н.).
        Но $\int X_n d\P = 1 \not\ra 0, n \ra +\infty$ -- следовательно, нет сходимости математических ожиданий.
        \item Заметим, что если из сходимости по вероятности следует сходимость в $L_1$, то так как из сходимости п.н. следует сходимость по $\P$, то из сходимости п.н. следует сходимость по $L_1$ -- что неверно по предыдущему пункту.
    \end{itemize}
\end{solution}